% !PRINTED 261
\weberpage{261}
%<<CONT
\weq{10}{\begin{aligned} A &= \left(\frac{u}{v}\right)^r + \left(\frac{v}{u}\right)^r \\ B &= (uv)^s + \left(\frac{2}{n}\right)^{r+s} \frac{2^s}{(uv)^s} , \end{aligned}}
in which $r$, $s$ are two integers that satisfy the conditions
\weq{11}{(n-1)r \equiv 0, \quad (n+1)s \equiv 0 \ (\mathrm{mod}\ 12)}
and which therefore, if $n$ is divisible by 3, are both divisible
by 3, and where
\[ \left(\frac{2}{n}\right) = (-1)^{\frac{n^2-1}{8}} \]
holds.

There then result the following mutually corresponding permutations:
\weq{12}{\begin{aligned}
&\omega, & & A, & & B, \\
&-\frac{1}{\omega}, & & A, & & B, \\
&\omega+2, & & (-1)^{\frac{(n-1)r}{12}} A, & & (-1)^{\frac{(n+1)s}{12}} B, \\
&\frac{\omega-1}{\omega+1}, & & \left(\frac{2}{n}\right)^r A, & & \left(\frac{2}{n}\right)^r B.
\end{aligned}}

If we now form a polynomial in $A$ and $B$ with numerical
coefficients
\weq{13}{\Phi_{a,c,\partial} = \Sigma\, M_{h,k}\, A^h B^k,}
in which, should changes of sign occur in (12), the exponents
$h$, $k$ are to be so arranged that in all terms of (13)
the same alterations of sign take place, then the system
of functions $\Phi_{a,c,\partial}$, or at least $\Phi^2_{a,c,\partial}$, will satisfy requirement 1.
of the theorem set up above; and in order to satisfy
requirement 2. as well, and so to obtain a transformation equation, we have to determine the coefficients $M_{h,k}$ in such a way
that all the $\nu$ values of $\Phi_{a,c,\partial}$ remain finite for $q=0$.
This problem is simplified essentially if $n$ has no square divisor, and still more if $n$ is a prime.

For if $n$ has no square divisor, one can derive all the
values $\Phi_{a,c,\partial}$ from $\Phi_{a,0,\partial}$ by increasing
$\omega$ by certain integers; if therefore, as is our requirement, no negative powers occur in the expansion of $\Phi_{a,0,\partial}$ % [sic] German source prints a period after 0 here; parallel occurrences have a comma.
in ascending powers
of $q$, then the same holds
of all the $\Phi_{a,c,\partial}$.
